\item En base a la evidencia clínica, se puede suponer que la tasa de disminución de la concentración $C$ de una droga en la sangre es proporcional a la propia concentración; esto es, $dC/dt = -kC$, en donde $k$ es una constante positiva. De modo que la concentración en el tiempo $t$ está dada por $C(t) = C_0 e^{-kt}$, en donde $C_0$ es la dosis inicial. Supóngase ahora que la droga se va a administrar en dosis iguales $C_0$ en intervalos de tiempo igualmente espaciados de longitud $T$. Supóngase que cuando se administra la droga se difunde inmediatamente por toda la corriente sanguínea; es decir, supóngase que hay una elevación instantánea de la concentración en cada administración de la droga.
    \begin{enumerate}
        \item Sea $C_{n-1}$ la concentración de droga inmediatamente después de la $n$-ésima inyección y sea $R_n$ la concentración inmediatamente antes de la $(n + 1)$-ésima inyección. De modo que $C_{n-1}$ y $R_n$ son las concentraciones al principio y final, respectivamente, del $n$-ésimo intervalo. Encuentre expresiones para $C_{n-1}$ y $R_n$.
        \item Calcule $\lim_{n \to \infty} C_{n-1}$ y $\lim_{n \to \infty} R_n$.
        \item Denótese por $C = C(t)$ la concentración de la droga en la sangre en cualquier instante $t$. Trace la gráfica de $C$.
        \item Supóngase que la droga es inefectiva debajo del nivel de concentración $C_L$ y que es nociva arriba del nivel de concentración $C_H$. Encuentre una dosis inicial $C_0$ y un intervalo de tiempo $T$ de manera que $C_L \le C(t) \le C_H$ para todo tiempo $t > 0$.
    \end{enumerate}
